\section{Fit a common scalar potential directly in the gradient domain}
Two independent approximations $\psi_d=f_d(i_d,i_q)$ and
$\psi_q=f_q(i_d,i_q)$ can have low sample error without correct parity or
integrability. Even parity-constrained independent fits do not guarantee a
common potential: $f_d=1+i_q^2$, $f_q=i_q$ have the required parities but
mixed derivatives $2i_q$ and zero. The missing information is the coupling
between their slopes.

The physically structured representation is
\begin{equation}
 \boxed{\hat\flux_{\rm fit}=\nabla\hat\coenergy,\qquad
 \hat\coenergy(i_d,i_q)=\hat\coenergy(i_d,-i_q).}\label{eq:fit-potential}
\end{equation}
A $C^2$ even scalar basis guarantees flux parity by the chain rule and
reciprocity by equality of mixed derivatives. Saturation and cross-saturation
are retained in the variation of its Hessian, not removed by symmetry.
The potential ties both flux surfaces to the same source of curvature.

Let $\vect c$ be its coefficients and $\vect\psi_k^{\rm rec}$ reconstructed
samples. The basic fitting problem is
\begin{equation}
 \boxed{\min_{\vect c}\sum_k
 \norm{\nabla\coenergy(\current_k;\vect c)-\vect\psi_k^{\rm rec}}^2
       +\lambda\mathcal R(\coenergy),\quad \coenergy(\current_0)=0.}
 \label{eq:gradient-fit}
\end{equation}
This asks for the admissible scalar landscape whose slopes best match the
observed vectors. It never requires directly measured energy samples.
If the terminal-current field is inconsistent with the declared magnetic
class, this minimization is a regularized projection onto that class under
the selected weights and basis. It does not establish a true storage law
in those terminal coordinates. Retain the sample residual
$\vect r_k=\vect\psi_k^{\rm rec}-\nabla\hat\coenergy(\current_k)$ alongside
predictions; constructed zero curl is not evidence that its discarded part
was entirely measurement error. Use
whitened residuals when uncertainties differ or share correlations; the
reconstruction divides voltage error by speed, so uniform flux uncertainty
should not be assumed across speeds.

Regularization controls degrees of freedom that measurements constrain weakly.
For example, a penalty on the spatial variation of the Hessian discourages
rapidly oscillating differential inductance. A penalty on the Hessian itself
also biases nonzero inductance toward zero; that distinction matters. Select
penalty strength using held-out regions and derivative targets, not training
error alone. Nondimensionalize current and flux so the numerical weight and
penalty have interpretable units. Robust losses may limit isolated outliers,
but cannot distinguish systematic bias from genuine magnetics.

Gauge removes the constant null direction. Additional rank deficiency can
remain if sample locations do not excite basis directions. Check the gradient
design rank on the gauge-fixed subspace and inspect scaled singular values
\cite{horn2013}. A smooth fitted potential is not necessarily locally convex:
positive differential inductance requires a separate Hessian constraint or
an explicit check over the supported domain.

\subsection{Add observation rows, then inspect their nullspace}
\label{sec:joint-gradient-fit}
The potential provides a common unknown field for electrical and mechanical
observations. Let $\matr G_k\vect c=\nabla W'(\vect i_k;\vect c)$
for a linear scalar basis, and let $\vect\eta$ collect explicitly modeled
measurement and loss parameters. A useful stacked residual is
\begin{equation}
 \vect r_k(\vect c,\vect\eta)=
 \begin{bmatrix}
 \vect u_k^{\rm meas}-R_s\vect i_k-\omega_{e,k}\matr J\matr G_k\vect c
                         -\vect d_{U,k}(\vect\eta)\\
 M_k^{\rm obs}-k_{\rm T}\vect a_k\trans\matr G_k\vect c
                         -d_{M,k}(\vect\eta)
 \end{bmatrix}.
 \label{eq:joint-observation-residual}
\end{equation}
Here the electrical discrepancy is in V and the torque discrepancy in Nm.
The observation $M^{\rm obs}$ needs a declared conversion to the modeled
torque, encoded in $d_M$ if justified. Unknown calibration parameters cannot
be hidden in arbitrary torque weights. With a justified joint covariance
$\matr\Sigma_k$, minimize
\begin{equation}
 \sum_k\vect r_k\trans\matr\Sigma_k^{-1}\vect r_k+
 \lambda\mathcal R(W'),\qquad W'(\vect i_0)=0 .
 \label{eq:joint-observation-fit}
\end{equation}
Whitening makes the observation part dimensionless and accommodates channel
correlation. Without calibrated uncertainty this is a formulation, not a
prescribed estimator with empirically chosen weights. If reconstructed flux
is used instead of voltage, its residual replaces the electrical block;
using both as independent data would count the same measurements twice.
For fixed nuisance parameters a torque row is
$k_{\rm T}\vect a_k\trans\matr G_k$. Reciprocity is already enforced by
the common potential, rather than added as another noisy observation.
Optional change penalties express a prior map preference; they do not
convert its torque-invisible components into measured ones.

Inspect the scaled Jacobian in $(\vect c,\vect\eta)$ on the gauge-fixed
subspace before claiming identification. A prior or smoothing penalty can
select a unique solution while the observation matrix still has a nullspace.
Conversely, a calibrated electrical block can remove radial modes that
torque plus reciprocity cannot remove. The supplied real channels below lack
the torque qualification needed to fit this joint magnetic model.

\paragraph{Keep fit quality separate from physical consistency.}
A small numerical residual measures agreement with measured samples. Parity,
curl, reciprocal mixed derivatives, differential-inductance smoothness, and
support coverage measure different properties. The admissible fit is the
best fit within the chosen conservative symmetric model, not automatically
the true machine. A larger structured residual can be more informative than
a nearly exact unconstrained interpolation of a distorted field. Excluded
physics, basis bias, and sensor errors can all contribute.
